\subsection{LINEAR MAPPINGS}

DEFINITION 4. Let E and F be two (left) modules with respect to the same ring A. A linear mapping (or A-linear mapping, or homomorphism, or A-homomorphism) of E into F is any mapping \(u: \mathbf{E} \to \mathbf{F}\) such that:

\begin{enumerate}
\def\labelenumi{(\arabic{enumi})}
\setcounter{enumi}{2}
\tightlist
\item
\end{enumerate}

\[
u (x + y) = u (x) + u (y) \quad for x \in \mathrm{E}, y \in \mathrm{E};\tag{4G}
\]

\[
u (\lambda . x) = \lambda . u (x) \quad for \lambda \in \mathrm{A}, x \in \mathrm{E}.
\]

If E and F are two right A-modules, a linear mapping \(u: \mathbf{E} \to \mathbf{F}\) is a mapping satisfying (3) and

\[
u (x . \lambda) = u (x). \lambda \quad \text { for } \lambda \in \mathrm{A}, x \in \mathrm{E}.\tag{4D}
\]

Remark. When E and F are two commutative groups considered as modules over the ring Z (no. 1), every homomorphism u of the group E (without operators) into the group F (without operators) is also a linear mapping of E into F, since for n an integer \textgreater0, the relation \(u(n.x) = n.u(x)\) follows from \(u(x + y) = u(x) + u(y)\) by induction on n and for n = -m \textless{} 0,

\[
u (n. x) = u (- (m. x)) = - u (m. x) = - (m. u (x)) = n. u (x).
\]

Examples. (1) Let E be an A-module and a an element of E; the mapping \(\lambda\mapsto\lambda a\) of the A-module \(A_{s}\) into E is a linear mapping \(\theta_{a}\) such that \(\theta_{a}(1)=a\) .

*(2) Let I be an open interval of the real line R, E the vector space of differentiable real-valued functions on I and F the vector space of all real-valued functions defined on I. The mapping \(x \mapsto x'\) which associates with every differentiable function x its derivative is a linear mapping from E to F.*

Note that a homothety \(x \mapsto \alpha x\) on an A-module E is not necessarily a linear mapping: in other words, the relation \(\alpha(\lambda x) = \lambda(\alpha x)\) does not necessarily hold for all \(\lambda \in A\) and \(x \in E\) . This relation is however true when \(\alpha\) belongs to the centre of A; \(x \mapsto \alpha x\) is then called a central homothety (cf.~no. 13).

If \(u: \mathbf{E} \to \mathbf{F}\) is a linear mapping, then, for every family \((x_{i})_{i \in I}\) of elements of \(\mathbf{E}\) and every family \((\lambda_{i})_{i \in I}\) of elements of \(\mathbf{A}\) such that the support of the family \((\lambda_{i}x_{i})_{i \in I}\) is finite,

\[
u \left(\sum_ {i \in I} \lambda_ {i} x _ {i}\right) = \sum_ {i \in I} \lambda_ {i} u (x _ {i})\tag{5}
\]

as follows immediately from (3) and (4G) by induction on the cardinal of the support of the family \((\lambda_{i}x_{i})\) .

PROPOSITION 1. Let E, F, G be three A-modules, u a linear mapping of E into F and v a linear mapping of F into G. Then the composite mapping \(v \circ u\) is linear.

PROPOSITION 2. Let E, F be two A-modules.

\begin{enumerate}
\def\labelenumi{(\arabic{enumi})}
\item
  If \(u: \mathbf{E} \to \mathbf{F}\) and \(v: \mathbf{F} \to \mathbf{E}\) are two linear mappings such that \(v \circ u\) is the identity mapping on \(\mathbf{E}\) and \(u \circ v\) is the identity mapping on \(\mathbf{F}\) , \(u\) is an isomorphism of \(\mathbf{E}\) onto \(\mathbf{F}\) and \(v\) is the inverse isomorphism.
\item
  Every bijective linear mapping \(u: \mathbf{E} \to \mathbf{F}\) is an isomorphism of \(\mathbf{E}\) onto \(\mathbf{F}\) .
\end{enumerate}

These propositions follow immediately from Definition 4.

Propositions 1 and 2 show that linear mappings can be taken as the morphisms for the species of A-module structures (Set Theory, IV, § 2, no. 1); we shall henceforth always assume that this choice of morphisms has been made.

Given two left (resp. right) A-modules E and F, let \(\text{Hom}(E, F)\) or \(\text{Hom}_{A}(E, F)\) denote the set of linear mappings of E into F.

The set \(\operatorname{Hom}(\mathbf{E},\mathbf{F})\) is a commutative ring, a subgroup of the product commutative group \(\mathbf{F}^{\mathbf{E}}\) of all mappings of \(\mathbf{E}\) into \(\mathbf{F}\) (I, § 4, no. 8); recall that for two elements \(u\) , \(v\) of \(\mathbf{F}^{\mathbf{E}}\) and for all \(x\in\mathbf{E}\) ,

\[
(u + v) (x) = u (x) + v (x), \quad (- u) (x) = - u (x)
\]

whence it follows immediately that, if u and v are linear, so are \(u + v\) and -u. If G is a third left (resp. right) A-module, \(f, f_{1}, f_{2}\) elements of Hom(E, F) and g, \(g_{1}, g_{2}\) elements of Hom(F, G), the following relations are immediately verified:

\begin{enumerate}
\def\labelenumi{(\arabic{enumi})}
\setcounter{enumi}{5}
\tightlist
\item
\end{enumerate}

\[
g \circ (f _ {1} + f _ {2}) = g \circ f _ {1} + g \circ f _ {2}\tag{7}
\]

\[
\left(g _ {1} + g _ {2}\right) \circ f = g _ {1} \circ f + g _ {2} \circ f\tag{8}
\]

\[
g \circ (- f) = (- g) \circ f = - (g \circ f).
\]

In particular, the law of composition \((f,g)\mapsto f\circ g\) on \(\operatorname{Hom}(\mathbf{E},\mathbf{E})\) defines with the above additive group structure a ring structure on \(\operatorname{Hom}(\mathbf{E},\mathbf{E})\) whose unit element, denoted by \(1_{E}\) or \(Id_{E}\) , is the identity mapping on E; the linear mappings of E into itself are also called endomorphisms of the A-module E and the ring \(\operatorname{Hom}(\mathbf{E},\mathbf{E})\) is also denoted by \(\operatorname{End}(\mathbf{E})\) or \(\operatorname{End}_{\mathbf{A}}(\mathbf{E})\) . The automorphisms of the A-module E are just the invertible elements of \(\operatorname{End}(\mathbf{E})\) (Proposition 2); they form a multiplicative group, denoted by \(\operatorname{Aut}(\mathbf{E})\) or \(\mathbf{GL}(\mathbf{E})\) which is also called the linear group relative to E.

It follows from (6) and (7) that, for two A-modules E, F, Hom(E, F) has the canonical structure of a left module over the ring Hom(F, F) and of a right module over Hom(E, E).

Let E, F, E', F' be four (left) A-modules, \(u: E' \to E\) and \(v: F \to F'\) A-linear mappings. If every element \(f \in \text{Hom}(E, F)\) is associated with the element \(v \circ f \circ u \in \text{Hom}(E', F')\) , a mapping

\[
\operatorname{Hom} (\mathbf {E}, \mathbf {F}) \rightarrow \operatorname{Hom} (\mathbf {E} ^ {\prime}, \mathbf {F} ^ {\prime})
\]

is defined, which is \(\mathbf{Z}\) -linear and which is denoted by \(\operatorname{Hom}(u,v)\) or \(\operatorname{Hom}_{\mathbf{A}}(u,v)\) . If \(u, u_1, u_2\) belong to \(\operatorname{Hom}(\mathbf{E}', \mathbf{E})\) and \(v, v_1, v_2\) to \(\operatorname{Hom}(\mathbf{F}, \mathbf{F}')\) , then

\[
\left\{ \begin{array}{l} \operatorname{Hom} (u _ {1} + u _ {2}, v) = \operatorname{Hom} (u _ {1}, v) + \operatorname{Hom} (u _ {2}, v) \\ \operatorname{Hom} (u, v _ {1} + v _ {2}) = \operatorname{Hom} (u, v _ {1}) + \operatorname{Hom} (u, v _ {2}) \end{array} \right.\tag{9}
\]

Let \(\mathbf{E}^{\prime \prime},\mathbf{F}^{\prime \prime}\) be two A-modules, \(u^{\prime}:\mathbf{E}^{\prime \prime}\to \mathbf{E}^{\prime}\) and \(v^{\prime}:\mathbf{F}^{\prime}\rightarrow \mathbf{F}^{\prime \prime}\) linear mappings. Then

\[
\operatorname{Hom} \left(u \circ u ^ {\prime}, v ^ {\prime} \circ v\right) = \operatorname{Hom} \left(u ^ {\prime}, v ^ {\prime}\right) \circ \operatorname{Hom} (u, v).\tag{10}
\]

If \(u\) is an isomorphism of \(\mathbf{E}'\) onto \(\mathbf{E}\) and \(v\) an isomorphism of \(\mathbf{F}\) onto \(\mathbf{F}'\) , \(\operatorname{Hom}(u, v)\) is an isomorphism of \(\operatorname{Hom}(\mathbf{E}, \mathbf{F})\) onto \(\operatorname{Hom}(\mathbf{E}', \mathbf{F}')\) whose inverse isomorphism is \(\operatorname{Hom}(u^{-1}, v^{-1})\) by (10).

If \(h\) (resp. \(k\) ) is an endomorphism of \(\mathbf{E}\) (resp. \(\mathbf{F}\) ), \(\operatorname{Hom}(h, 1_{\mathbf{F}})\) (resp. \(\operatorname{Hom}(1_{\mathbf{E}}, k)\) ) is just the homothety with ratio \(h\) (resp. \(k\) ) for the right (resp. left) module structure on the ring \(\operatorname{End}(\mathbf{E})\) (resp. \(\operatorname{End}(\mathbf{F})\) ) defined above.

